\TopicDivider{01}{Epistemic Uncertainty \& Bayes' Rule}{Probability as a calculus of rational belief under incomplete information.}

\begin{frame}{First-Principles Epistemology}
  \begin{LearningGoals}{Foundational Learning Invariants}
    \begin{itemize}\setlength{\itemsep}{.35em}
      \item Differentiate \textbf{epistemic uncertainty} (reducible via data) from \textbf{aleatoric variability} (irreducible stochasticity).
      \item Formulate Bayes' Theorem as an exact consequence of Kolmogorov's conditional probability axioms.
      \item Diagnose why the marginal likelihood denominator \(\int_{\Theta} p(y \mid \theta) p(\theta) \, d\theta\) forms the primary computational challenge in machine learning.
    \end{itemize}
  \end{LearningGoals}
  \vspace{.15cm}
  \TakeawayBanner{Probability in Bayesian inference is a measure of belief under incomplete information, not asymptotic long-run frequency.}
\end{frame}

\begin{frame}{The Triad of Bayesian Components}
  \StatGridThree{%
    \StatCard{$P(\theta)$}{Prior Space}{Initial Hypothesis Density}{Encodes domain physics \& boundary constraints}%
  }{%
    \StatCard{$P(y \mid \theta)$}{Likelihood}{Generative Plausibility}{Forward observational process model}%
  }{%
    \StatCard{$P(\theta \mid y)$}{Posterior Space}{Rational Belief Update}{Conditioned on observed empirical evidence}%
  }
  \vspace{.15cm}
  \begin{ConceptCard}[INVARIANT]{The Conditioning Principle}
    Observing realization \(y\) filters out impossible states of nature, renormalizing probability mass exclusively across surviving hypotheses consistent with data.
  \end{ConceptCard}
\end{frame}

\begin{frame}{The Anatomy of Coherent Belief Revision}
  \TwoPane{%
    \begin{ConceptCard}[FOUNDATIONS]{The Generative Invariant}
      In Bayesian probability theory, parameters are not fixed unknown constants with sampling variance; they are random variables with explicit probability distributions.
      \vspace{.25em}
      \begin{itemize}\setlength{\itemsep}{.2em}
        \item \textbf{Prior} \(p(\theta)\): Hypotheses distribution before observations.
        \item \textbf{Likelihood} \(p(y \mid \theta)\): Conditional probability of data given state.
      \end{itemize}
    \end{ConceptCard}
  }{%
    \begin{DeepDiveCard}[THE BOTTLENECK]{The Evidence Integral}
      The marginal likelihood \(p(y)\) normalizes the posterior:
      \begin{equation*}
        p(y) = \int_{\Theta} p(y \mid \theta) \, p(\theta) \, d\theta
      \end{equation*}
      In \(D\) dimensions (\(D > 50\)), grid integration requires \(O(K^D)\) evaluations—an exponential catastrophe necessitating MCMC and VI.
    \end{DeepDiveCard}
  }
\end{frame}

\begin{frame}{Exact Formulation of Bayes' Rule}
  \MathDerivationStep{1}{Theorem Statement}{p(\theta \mid y) = \frac{p(y \mid \theta) \, p(\theta)}{\int_{\Theta} p(y \mid \theta') \, p(\theta') \, d\theta'}}{%
    The joint probability \(p(y, \theta)\) factors as \(p(y \mid \theta)p(\theta)\) and symmetrically as \(p(\theta \mid y)p(y)\). Equating the joint representations yields the theorem directly.
  }
  \vspace{.12cm}
  \MathDerivationStep{2}{Proportional Kernel}{p(\theta \mid y) \propto p(y \mid \theta) \times p(\theta)}{%
    Because \(p(y)\) is independent of \(\theta\), optimization and sampling (e.g., MCMC, HMC) depend strictly on the unnormalized posterior kernel.
  }
\end{frame}

\begin{frame}{Visualizing Epistemic Variance Collapse}
  \begin{ScientificFigure}{PGFPlots Density Shift Under Evidence Accumulation}
    \centering
    \begin{tikzpicture}
      \begin{axis}[bayes/science plot,xlabel={Parameter Space \(\theta \in [0, 10]\)},ylabel={Probability Density \(p(\theta)\)},xmin=0,xmax=10,ymin=0,ymax=1.2,legend pos=north west]
        \addplot[color=BayesMuted,line width=1.5pt,domain=0:10,samples=100] {exp(-0.15*(x-3.0)^2)};
        \addlegendentry{Diffuse Prior: \(\mathcal{N}(3.0, 3.3)\)}
        
        \addplot[color=BayesAmber,line width=1.5pt,domain=0:10,samples=100] {exp(-0.45*(x-6.5)^2)};
        \addlegendentry{Likelihood Profile: \(L(\theta \mid y)\)}
        
        \addplot[color=BayesTeal,line width=2.0pt,domain=0:10,samples=100] {1.12*exp(-0.85*(x-5.9)^2)};
        \addlegendentry{Posterior Distribution: \(p(\theta \mid y)\)}
      \end{axis}
    \end{tikzpicture}
    \BayesCaption{Observing data concentrates probability mass, contracting variance from the diffuse prior to the sharp posterior.}
  \end{ScientificFigure}
\end{frame}
